% Part: incompleteness
% Chapter: incompleteness-provability
% Section: introduction

\documentclass[../../../include/open-logic-section]{subfiles}

\begin{document}

% SEGMENT OLP-0313-S01

\olfileid{inc}{inp}{int}

\olsection{முன்னுரை}

இயல் எண்கள் போன்ற ஒரு கணிதக் கட்டமைப்புக்கான அடிகோள்களின் முறை,
அக்கட்டமைப்பைப் பற்றிய எல்லா மெய்யான கூற்றுகளையும் வருவிக்க அனுமதிக்காவிட்டால்
போதுமானதல்ல என்று ஹில்பர்ட் கருதினார். வருவித்தலின் முறைசார் அமைப்புகளில்
அவருக்கு பின்னர் ஏற்பட்ட ஆர்வத்துடன் இதை இணைத்துப் பார்த்தால், இயல் எண்களைப்
பற்றி நாம் காரணங்கூறும் முறைசார் அமைப்புகள் முரணற்றவை மட்டுமல்ல,
\emph{முழுமையானவையும்} ஆக வேண்டும் என்று அவர் கருதியிருப்பார்: அதாவது, அதன்
மொழியிலுள்ள ஒவ்வொரு கூற்றும் !!{derivable} ஆக இருக்க வேண்டும், அல்லது
அதன் மறுதலை வருவிக்க முடியும். கோடலின் முதல்
முழுமையின்மைத் தேற்றம் அத்தகைய அடிகோள்முறை எதுவும் இல்லை என்பதைக்
காட்டுகிறது: எண்கணிதத்திற்கான முழுமையான, முரணற்ற, !!{axiomatizable}
முறைசார் கோட்பாடு எதுவும் இல்லை. உண்மையில், எந்த “போதிய வலிமை”
கொண்ட முரணற்ற, !!{axiomatizable} கணிதக் கோட்பாடும் முழுமையானதல்ல.

ஹில்பர்ட்டின் இன்னும் முக்கியமான குறிக்கோள், நவீன (“செவ்வியல்”)
கணிதத்தை நியாயப்படுத்தும் அவரது திட்டத்தின் மையக்கருத்தாக இருந்தது:
செவ்வியல் காரணங்கூறலைப் பிரதிநிதித்துவப்படுத்தும் முறைசார் அமைப்புகளுக்கான
வரையறுக்கப்பட்ட எண்ணியல் முரணின்மை நிறுவல்களைக் கண்டுபிடிப்பது. ஆகவே
ஹில்பர்ட்டின் திட்டத்தைப் பொறுத்தவரை, கோடலின் இரண்டாம் முழுமையின்மைத்
தேற்றம் மிகப் பெரிய பின்னடைவாக இருந்தது. முதல் தேற்றத்தைப் போலவே,
இரண்டாம் தேற்றத்தையும் தெளிவற்ற சொற்களில் கூறலாம். தோராயமாகச் சொன்னால்,
எண்கணிதத்தின் எந்தப் போதிய வலிமை கொண்ட கோட்பாடும் தன் சொந்த
முரணின்மையை நிறுவ முடியாது என்று அது கூறுகிறது. இங்கு “போதிய
வலிமை” என்பது $\Th{Q}$ ஐவிட சற்றுக் கூடுதலான வலிமையை உள்ளடக்கியதாக
இருக்க வேண்டும்.

முழுமையின்மைத் தேற்றத்திற்கான கோடலின் முதற் நிறுவலின் எண்ணத்தை
எபிமெனிடீஸின் முரண்பாட்டில் காணலாம். கிரீட் நாட்டைச் சேர்ந்த எபிமெனிடீஸ்,
அனைத்து கிரீட்டியர்களும் பொய்யர்கள் என்று கூறினார்; முரண்பாட்டின் இன்னும்
நேரடியான வடிவம் “இந்த வாக்கியம் பொய்யானது” என்பதாகும். மெய்மைக்குப்
பதிலாக !!{derivability} ஐ வைத்ததன் மூலம், கோடல் தன்னையே
!!{derivable} அல்ல என்று சுற்றிவளைத்துக் கூறும் !!a{sentence} ஒன்றை
முறைப்படுத்த முடிந்தது. அந்த !!{sentence} !!{derivable} ஆக இருந்தால்,
கோட்பாடு முரணாகிவிடும். தாம் கருதிய அடிகோள்முறையிலிருந்து அந்த
!!{sentence} இன் மறுதலையும் !!{derivable} அல்ல என்பதை கோடல் காட்டினார்.
இந்த இரண்டாம் பகுதிக்காக, கோட்பாடு $\Th{T}$ “$\omega$-முரணற்றது”
என்று அழைக்கப்படுவதை அவர் கருத வேண்டியிருந்தது.
$\omega$-முரணற்ற தன்மை முரணற்ற தன்மையுடன் தொடர்புடையது, ஆனால் அது
அதைவிட வலிமையான பண்பு.
\footnote{அதாவது, ஒவ்வொரு $\omega$-முரணற்ற கோட்பாடும் முரணற்றது;
ஆனால் இதன் மறுதிசை பொதுவாக உண்மையல்ல.}
கோடலுக்குப் சில ஆண்டுகளுக்குப் பின்னர், $\Th{T}$ இன் எளிய முரணற்ற
தன்மையை மட்டும் கருதினால் போதும் என்று ரோஸர் காட்டினார்.

தன்னையே குறிப்பிடும் !!a{sentence} ஒன்றை எவ்வாறு உருவாக்குவது என்பதே
முதல் சவால். $\Th{Q}$ இன் மொழியிலுள்ள ஒவ்வொரு !!{formula}~$!A$ க்கும்,
$\Gn{!A}$ க்கு ஒத்த எண்ணுருவை $\gn{!A}$ என்று குறிக்கலாம். இதன் பொருளை
கவனியுங்கள்: $!A$ என்பது $\Th{Q}$ இன் மொழியிலுள்ள !!a{formula};
$\Gn{!A}$ ஓர் இயல் எண்; $\gn{!A}$ என்பது $\Th{Q}$ இன் மொழியிலுள்ள
\emph{சொல்}. எனவே $\Th{Q}$ இன் மொழியிலுள்ள ஒவ்வொரு !!{formula}~$!A$
க்கும் $\gn{!A}$ என்ற \emph{பெயர்} உண்டு; அது $\Th{Q}$ இன் மொழியிலுள்ள ஒரு
சொல். இதன் மூலம் $\Th{Q}$ இன் மொழியிலுள்ள !!{formula}s பிற
!!{formula}s பற்றிக் “கூற” முடியும் என்ற கருத்துக் கட்டமைப்பு
கிடைக்கிறது. பின்வரும் துணைத்தேற்றம் நிலைப்புள்ளித் துணைத்தேற்றம்
என்று அழைக்கப்படுகிறது.

\begin{lem}
எந்த $\Th{Q}$ நீட்சியாக இருந்தாலும் $\Th{T}$ ஐ எடுத்துக்கொள்க; மேலும்
$x$ மட்டுமே கட்டற்ற மாறியாக உள்ள ஏதேனும் !!{formula}~$!B(x)$ ஐ
எடுத்துக்கொள்க. அப்போது
$\Th{T} \Proves!A \liff !B(\gn{!A})$ ஆகுமாறு !!a{sentence}~$!A$
ஒன்று உள்ளது.
\end{lem}

எந்தப் பண்பு $!B(x)$ கொடுக்கப்பட்டாலும், “$!B(x)$ எனக்கு உண்மை”
என்று கூறும் !!a{sentence}~$!A$ ஒன்று உள்ளது, மேலும் $\Th{T}$ இதை
“அறியும்” என்பதே இந்தத் துணைத்தேற்றத்தின் கூற்று.

இத்தகைய !!a{sentence} ஐ எவ்வாறு உருவாக்கலாம்? குவைன் முன்வைத்த
எபிமெனிடீஸ் முரண்பாட்டின் பின்வரும் வடிவத்தைக் கருதுக:
\begin{quote}
“தன் மேற்கோளால் முன்னிடப்பட்டபோது பொய்யைத் தருகிறது” என்பது தன்
மேற்கோளால் முன்னிடப்பட்டபோது பொய்யைத் தருகிறது.
\end{quote}
இந்த வாக்கியம் நேரடியாகத் தன்னைக் குறிப்பிடுவதில்லை. மேற்கோள் குறிகளுக்குள்
உள்ள தொடரியல் பொருட்களைப் பற்றி மட்டுமே அது ஒரு கூற்றைச் செய்கிறது;
இந்த வகையில் அது பின்வரும் வாக்கியங்களைப் போன்றதே:
\begin{enumerate}
\item “ராபர்ட்” என்பது ஒரு நல்ல பெயர்.
\item “நான் ஓடினேன்.” என்பது ஒரு குறுகிய வாக்கியம்.
\item “மூன்று சொற்களைக் கொண்டது” மூன்று சொற்களைக் கொண்டது.
\end{enumerate}
ஆனால் “தன் மேற்கோளால் முன்னிடப்பட்டபோது பொய்யைத் தருகிறது” என்ற
சொற்றொடரை எடுத்து, அதற்கே மேற்கோள் வடிவில் தன்னையே முன்னிட்டால் என்ன
நடக்கும்? அப்போது தொடக்க வாக்கியமே கிடைக்கிறது!{} சுருக்கமாக, அந்த
வாக்கியம் தானே பொய்யானது என்று கூறுகிறது.

\end{document}
